\section{Lecture 19: Analytic Stacks --- Definition, Examples, and the Tate Elliptic Curve}
\label{lec:19}

Having spent \cref{lec:16,lec:17,lec:18} building the $!$-topology on
$\Aff=(\AnRing)^{\mathrm{op}}$, Scholze now finally defines an
\emph{analytic stack} --- and, as promised, it is not difficult: a sheaf on
$\Aff$ for the $!$-topology, packaged as an accessible functor
$\AnRing\to\operatorname{Anima}$. The bulk of the lecture is a tour of examples,
showing that essentially every existing flavour of geometry --- derived schemes,
adic spaces, complex- and real-analytic spaces, topological/$C^k$/$C^\infty$
manifolds, condensed anima --- embeds into analytic stacks, often in several
inequivalent ways related by natural comparison maps (one incarnation of GAGA
being an isomorphism of stacks). The lecture closes by returning to the Tate
elliptic curve $E_q$, now constructible over the gaseous base ring via a
``norm'' morphism $\mathbb P^1\to[0,\infty]$.

Throughout, ``ring'' means commutative ring and everything is (light) condensed.

\subsection{The definition}
\label{subsec:19:definition}

Recall (\cref{lec:15}) that $\AnRing$, the $\infty$-category of analytic rings,
is presentable: Clausen showed it has all colimits and is generated under
colimits by a set's worth of objects (the $\kappa$-compact analytic rings for
$\kappa=(2^{\aleph_0})^+$, \cref{thm:15:anring-presentable}). This is what makes
the following definition sensible.

\begin{definition}[Analytic stack]
\label{def:19:analytic-stack}
An \emph{analytic stack} is an accessible functor
\begin{equation*}
X\colon\ \AnRing\ \longrightarrow\ \operatorname{Anima}
\end{equation*}
that commutes with finite products and satisfies the following descent
condition: for every $!$-hypercover $\AnSpec(A_\bullet)\to\AnSpec(A)$
\emph{for which} the associated diagram of derived categories is a limit,
\begin{equation}
\label{eq:19:derived-descent}
D(A)\ \xrightarrow{\ \sim\ }\ \varprojlim_{\Delta}{}^{!}\,D(A_\bullet),
\end{equation}
the space of $A$-points is the corresponding limit of the $A_\bullet$-points:
\begin{equation}
\label{eq:19:stack-descent}
X(A)\ \xrightarrow{\ \sim\ }\ \varprojlim_{\Delta}\,X(A_\bullet).
\end{equation}
Here \emph{accessible} means $X$ commutes with $\kappa$-filtered colimits for
some (equivalently, all sufficiently large) regular cardinal $\kappa$; the word
is present only to circumvent the set-theoretic issue that $\AnRing$ is not
small (\cref{subsec:15:set-theory}). Commuting with finite products (including
the empty product) says geometrically that $X$ turns finite disjoint unions of
affines into products. We write $\AnStk$ for the $\infty$-category of analytic
stacks.
\end{definition}

The subtlety is that descent is imposed only along the hypercovers that already
satisfy the derived condition \eqref{eq:19:derived-descent}. Recall
(\cref{lec:17,cor:17:descent-codescent}) that \eqref{eq:19:derived-descent} is
equivalently the statement that $D(A)$ is the colimit of the $D(A_\bullet)$
along the lower-shriek functors, computed in $\PrL$:
\begin{equation*}
D(A)\ \xrightarrow{\ \sim\ }\ \varinjlim_{\Delta^{\mathrm{op}}}{}^{!}\,D(A_\bullet)
\quad\text{in }\PrL,
\end{equation*}
a condition stable under base change.

\begin{remark}[An intermediate topos]
\label{rmk:19:intermediate-topos}
The definition selects, up to set theory, an $\infty$-topos lying strictly
\emph{between} $!$-sheaves and $!$-hypersheaves. For \v Cech nerves of
$!$-covers the condition \eqref{eq:19:derived-descent} holds automatically (by
the very definition of a $!$-cover, \cref{def:18:shriek-cover}), so every
analytic stack is a $!$-sheaf. For hypercovers one would normally demand descent
against \emph{all} hypercovers; here one asks it only against those hypercovers
for which the derived categories already glue, i.e.\ satisfy
\eqref{eq:19:derived-descent}. Restricting the class of hypercovers this way is
exactly what is needed so that the primary invariant $D(X)$ (below) will itself
satisfy descent.
\end{remark}

\begin{remark}[Colimits and sheafification]
\label{rmk:19:sheafification}
One can form colimits in $\AnStk$: take the colimit in accessible presheaves
$\operatorname{Fun}(\AnRing,\operatorname{Anima})$ --- still an accessible functor
--- and then enforce the sheaf condition by sheafifying. Sheafification preserves
accessibility (and set-theoretic size); this is the analogue for the
$!$-topology of Waterhouse's theorem \cite{Waterhouse_1975} for the fpqc
topology (\cref{rmk:15:accessible-representable}). Constructing this
sheafification requires knowing that the pullback of a hypercover satisfying
\eqref{eq:19:derived-descent} again satisfies it, which follows from the same
base-change argument as in \cref{lec:17,lec:18}.
\end{remark}

\subsection{Examples}
\label{subsec:19:examples}

\begin{example}[Affine analytic stacks]
\label{ex:19:affine}
For any $A\in\AnRing$ the functor of points
\begin{equation*}
\AnSpec(A)\colon\ B\ \longmapsto\ \Homop_{\AnRing}(A,B)
\end{equation*}
is an analytic stack. It manifestly commutes with finite products (it lands in
mapping \emph{spaces}, already in $\operatorname{Anima}$, so nothing needs
sheafifying), and the descent condition \eqref{eq:19:stack-descent} is precisely
that $A$ is recovered as the limit of the $A_\bullet$ --- both on the level of
rings and, using \eqref{eq:19:derived-descent}, on the level of derived
categories. Thus one obtains a fully faithful embedding
\begin{equation}
\label{eq:19:anring-embed}
(\AnRing)^{\mathrm{op}}=\Aff\ \hookrightarrow\ \AnStk .
\end{equation}
The affine analytic stacks are the basic building blocks; all others arise by
gluing them, i.e.\ as colimits.
\end{example}

\begin{remark}[Reading the accessibility condition]
\label{rmk:19:accessible-colimit}
The accessibility clause in \cref{def:19:analytic-stack} is exactly the
condition that $X$ be a small colimit of objects in the essential image of
\eqref{eq:19:anring-embed}:
\begin{equation*}
X\ \text{accessible}\ \iff\ X\ \text{is a small colimit of affines }\AnSpec(A).
\end{equation*}
The colimit here is taken in presheaves (before enforcing the sheaf condition);
sufficiently filtered colimits already preserve the descent condition, so one may
sheafify afterwards. This is the same dichotomy as for condensed sets
(\cref{lec:2}), which one may think of either as built from profinite sets under
small colimits or as the accessible sheaves.
\end{remark}

The point of imposing so strong a topology is that analytic stacks with
completely different presentations --- built as colimits of affines in a priori
unrelated ways --- can nonetheless be identified. A strong topology is what makes
such non-obvious isomorphisms hold; we will see a GAGA instance below.

\begin{example}[Derived schemes]
\label{ex:19:derived-schemes}
There is a functor
\begin{equation}
\label{eq:19:dersch-to-anstk}
\operatorname{DerSch}\ \longrightarrow\ \AnStk,\qquad
\operatorname{Spec}(A)\ \longmapsto\ \AnSpec(A),
\end{equation}
where a derived scheme is regarded as a sheaf for the descendable topology on
affine derived schemes and $A$ is given the \emph{trivial} analytic ring
structure (all condensed $A$-modules complete). Since descendable covers go to
$!$-covers (\cref{prop:18:zariski,prop:18:descendable}), this is well-defined. It is fully
faithful on derived schemes; the proof uses Bhatt's Tannaka duality.%
\footnote{Scholze named ``Bhatt's Tannaka duality'' without a precise reference.
Editorially, the relevant statement is Bhatt's \emph{Algebraization and Tannaka
duality} \cite{bhatt2014algebraizationtannakaduality}, where a
qcqs (derived) scheme or algebraic space $X$ is recovered from the symmetric
monoidal $\infty$-category $\mathrm{QCoh}(X)$, so that maps into $X$ correspond
to suitable tensor functors on quasi-coherent complexes.}
\end{example}

\begin{example}[Adic spaces]
\label{ex:19:adic}
Assuming sheafiness throughout (\cref{def:10:sheafy}), there is a functor
\begin{equation}
\label{eq:19:adic-to-anstk}
\{\text{adic spaces}\}\ \longrightarrow\ \AnStk,\qquad
\Spa(A,A^+)\ \longmapsto\ \AnSpec\bigl((A,A^+)_\solid\bigr).
\end{equation}
It is fully faithful on quasi-projective objects; one does not know a
full faithfulness statement in general, because there is no good version of
Tannaka duality for adic spaces.
\end{example}

\begin{remark}[Two embeddings of schemes into adic/analytic spaces]
\label{rmk:19:two-embeddings}
Composing $\operatorname{Sch}\hookrightarrow\{\text{adic spaces}\}$ with
\eqref{eq:19:adic-to-anstk} gives \emph{different} functors, according to the
choice of ring of integral elements. For a discrete ring $A$,
\begin{equation*}
\begin{tikzcd}[row sep=small, column sep=large]
& \Spa(A,\mathbb Z)\arrow[r,mapsto] & \AnSpec\bigl((A,\mathbb Z)_\solid\bigr) \\
\operatorname{Spec}(A)\arrow[ur,mapsto]\arrow[dr,mapsto] & & \\
& \Spa(A,A)\arrow[r,mapsto] & \AnSpec\bigl((A,A)_\solid\bigr).
\end{tikzcd}
\end{equation*}
Both are fully faithful embeddings. They differ sharply in their geometry, as
already flagged by Clausen (\cref{rmk:18:fuzz}): under the first, $(A,\mathbb
Z)_\solid$, \emph{every} map of affines becomes proper and open immersions become
closed immersions (Zariski opens ``look closed''); under the second,
$(A,A)_\solid$ (the ``relatively solid'' structure), open immersions go to open
immersions, matching the usual scheme-theoretic notions. The topmost row
$(A,\mathbb Z)_\solid$ is itself obtained from the basic embedding into
$\AnSpec(\mathbb Z_\solid)$ by base change from $\AnSpec(\mathbb Z)$ with the
trivial structure. More generally, for any $A\in\AnRing$ one gets a functor
\begin{equation*}
\operatorname{DerSch}_{\ur A(\ast)}\ \longrightarrow\ \AnStk_{/A},
\end{equation*}
from derived schemes over the underlying discrete ring $\ur A(\ast)$ to analytic
stacks over $A$, by using the previous functors and base-changing along
$\AnSpec(A)\to\AnSpec(\ur A(\ast))$. For instance, taking $A=\mathbb C$ with its
gaseous structure, ordinary schemes over $\mathbb C$ map to analytic stacks over
$\mathbb C_{\mathrm{gas}}$.
\end{remark}

\begin{example}[Complex-analytic spaces]
\label{ex:19:complex-analytic}
There is a functor
\begin{equation*}
\{\text{complex-analytic spaces}\}\ \longrightarrow\ \AnStk_{/\mathbb C_{\mathrm{gas}}},
\qquad
X=\bigcup_K K\ \longmapsto\ \bigcup_{K\subseteq X}\AnSpec\,\mathcal O(K)^{t},
\end{equation*}
where $K$ ranges over the \emph{compact Stein} subsets of $X$. Complex-analytic
geometers do not usually isolate a notion of ``affinoid'', but they could: any
complex-analytic space is a union of compact Stein subsets, and these behave very
much like affinoid objects. A typical compact Stein is the vanishing locus of an
ideal in a closed polydisc,
\begin{equation}
\label{eq:19:compact-stein}
V(I)\subseteq\mathbb D^n=\{(z_1,\dots,z_n)\in\mathbb C^n:\ |z_i|\le1\}.
\end{equation}
On such a $K$ one puts the algebra of \emph{overconvergent} holomorphic functions
\begin{equation}
\label{eq:19:overconvergent}
\mathcal O(K)^{t}\ =\ \varinjlim_{U\supseteq K}\mathcal O(U),
\end{equation}
the colimit over open neighbourhoods $U$ of $K$ of the holomorphic functions on
$U$; it carries a natural (dual nuclear Fréchet, DNF) topology, hence a condensed
ring structure that is a gaseous $\mathbb C$-algebra, and one gives $\AnSpec$ the
induced analytic ring structure. The functor is fully faithful on quasi-projective
objects.
\end{example}

\begin{remark}[Noetherianity of $\mathcal O(K)$]
\label{rmk:19:stein-noetherian}
It is a theorem of the 1970s--80s that for $K$ suitably close to a polydisc ---
in particular for $K=V(I)$ a Zariski-closed subset of a polydisc as in
\eqref{eq:19:compact-stein} --- the ring $\mathcal O(K)^t$ is Noetherian and
excellent, and regular when $X$ is a manifold, so that coherent sheaves on $X$ are
described locally by finitely generated modules over these Noetherian rings.
There is a subtlety: for a more general compact Stein (e.g.\ one whose slice is a
profinite subset of the polydisc, giving infinitely many connected components of a
zero locus) Noetherianity can fail; it is the Zariski-closed-in-a-polydisc case
that is Noetherian.%
\footnote{Scholze gave no precise attribution. Editorially: Noetherianity of
$\mathcal O(K)$ for $K$ a compact Stein set is due to Frisch
\cite{Frisch_1967}, and excellence to Bingener \cite{Bingener_1978} and
Scheja--Storch \cite{Scheja_1972}; cf.\ the appendix of Clausen--Scholze's
lectures on condensed mathematics and complex geometry.}
\end{remark}

\begin{remark}[Other manifold flavours]
\label{rmk:19:manifold-flavours}
The same kind of definition works for real-analytic manifolds, smooth ($C^\infty$)
manifolds, $C^k$-manifolds, and $C^0=$ topological manifolds; one uses in each
case the appropriate over-convergent functions on neighbourhoods (over-convergent
continuous functions being a slightly awkward but workable notion). In each case
one may take real- or complex-valued functions, giving two functors with one the
base change of the other along $\mathbb R\to\mathbb C$.
\end{remark}

\begin{example}[Condensed anima]
\label{ex:19:condensed-anima}
Light condensed anima also map to analytic stacks:
\begin{equation*}
\operatorname{CondAn}^{\mathrm{light}}\ \longrightarrow\ \AnStk .
\end{equation*}
This comes not from the condensed structure on the rings but from the test
objects: a light profinite set $S\in\ProFin_{\mathbb N}(\mathrm{Fin})$ is sent to
the spectrum of its ring of continuous ($=$ locally constant) $\mathbb Z$-valued
functions,
\begin{equation*}
S\ \longmapsto\ \operatorname{Spec}\,\Cont(S,\mathbb Z)\ \longmapsto\ \AnSpec ,
\end{equation*}
composed with the derived-scheme functor \eqref{eq:19:dersch-to-anstk} (with its
trivial structure). For $S$ finite this is a finite disjoint union of copies of
$\operatorname{Spec}\mathbb Z$; in general $S=\varprojlim_i S_i$ and the spectrum
is the corresponding limit. A surjection $S'\twoheadrightarrow S$ induces a
faithfully flat map $\Cont(S,\mathbb Z)\to\Cont(S',\mathbb Z)$, and
\emph{lightness} guarantees it is countably presented, hence descendable
(\cref{ex:18:descendable}), so a $!$-cover; this is why condensed anima, defined
as hypersheaves of anima on $\ProFin_{\mathbb N}(\mathrm{Fin})$, map to analytic
stacks. The intermediate topology between sheaves and hypersheaves reappears here:
one checks that a hypercover of light profinite sets goes to something satisfying
exactly the descent condition \eqref{eq:19:derived-descent}.
\end{example}

\subsection{Incarnations and a GAGA isomorphism}
\label{subsec:19:incarnations}

The examples above give \emph{several} incarnations of the same underlying object
as an analytic stack, related by natural maps. Take the two-sphere $S^2$, worked
over the gaseous complex numbers throughout.

\begin{remark}[Incarnations of $S^2\cong\mathbb P^1_{\mathbb C}$]
\label{rmk:19:incarnations}
There are natural maps
\begin{equation}
\label{eq:19:incarnations}
\begin{tikzcd}[row sep=large, column sep=small]
(S^2)^{\mathrm{top.mfd}}_{\mathbb C_{\mathrm{gas}}}=\AnSpec\,\Cont(S^2,\mathbb C)_{\mathrm{gas}}
\arrow[d] \\
(S^2)^{\mathbb R\text{-}\mathrm{an}}_{\mathbb C_{\mathrm{gas}}}=\AnSpec\,C^{\omega}(S^2,\mathbb C)_{\mathrm{gas}}
\arrow[d] \\
(\mathbb P^1_{\mathbb C})^{\mathbb C\text{-}\mathrm{an}}_{\mathbb C_{\mathrm{gas}}}
\arrow[d, "\sim"', "\mathrm{GAGA}"] \\
(\mathbb P^1_{\mathbb C})^{\mathrm{alg}}_{\mathbb C_{\mathrm{gas}}}\, .
\end{tikzcd}
\end{equation}
The top two are affine: $S^2$ as a topological manifold gives
$\AnSpec$ of the (gaseous) algebra of continuous functions, and as a
real-analytic manifold gives $\AnSpec$ of the algebra $C^\omega(S^2,\mathbb C)$
of real-analytic functions. The complex-analytic $\mathbb P^1$ is \emph{not}
affine: it is glued from two copies of $\AnSpec\,\mathcal O(\mathbb D)^t$
(over-convergent holomorphic functions on a disc) along the over-convergent
holomorphic functions on the circle $S^1$. The algebraic $\mathbb P^1$ is
likewise glued from two copies of $\AnSpec$ of the polynomial algebra
$\mathbb C[T]$. In addition, one may treat $S^2$ purely as a condensed set and
base-change to $\mathbb C_{\mathrm{gas}}$ (\cref{ex:19:condensed-anima}); this
incarnation is secretly defined over $\mathbb Z$ (it uses locally constant
$\mathbb Z$-valued functions on a profinite presentation
$S^2=\widetilde S/\!\sim$), and there is a natural map from the continuous-function
incarnation at the top all the way down to it.
\end{remark}

\begin{remark}[GAGA as an isomorphism of stacks]
\label{rmk:19:gaga}
The marked arrow in \eqref{eq:19:incarnations} is an \emph{isomorphism}: the
complex-analytic $\mathbb P^1$ and the algebraic $\mathbb P^1$ agree as analytic
stacks. This is one incarnation of GAGA --- and a strong one, an isomorphism of
the geometric objects themselves (in particular of their derived categories), not
merely a comparison of coherent sheaves or their cohomology. The two sides are
built in genuinely different ways --- one from polynomial algebras, one from
algebras of over-convergent holomorphic functions --- and it is precisely the
strong ($!$-)topology imposed in \cref{def:19:analytic-stack} that forces such an
isomorphism to hold. One proves it by covering both by two discs, checking the
two discs form a descendable ($!$-)cover, and unravelling what it means for the
glued objects to coincide. (The GAGA \emph{isomorphism} of stacks is logically
prior to, and does not by itself contain, statements like Riemann--Roch, which
require knowing the object is algebraic.)
\end{remark}

\begin{remark}[The sheaf $X\mapsto D(X)$ and $\PrL$-valued sheaves]
\label{rmk:19:D-sheaf}
The assignments
\begin{equation*}
\AnSpec(A)\ \longmapsto\ D(A)\ \longmapsto\ \PrL_{D(A)}=\Mod_{D(A)}(\PrL)\ \in\ \Catinf
\end{equation*}
satisfy descent (this was the design goal of the $!$-topology), so they induce,
for every analytic stack $X\in\AnStk$, an $\infty$-category $D(X)$ of
quasi-coherent sheaves and a $2$-category $\PrL_X$ of sheaves of categories over
$X$, via the intermediate (sheaf-but-not-quite-hypersheaf) descent condition.
Both are part of a six-functor formalism to be developed later. The different
incarnations of \cref{rmk:19:incarnations} have very different module theories:
for the top (continuous-function) incarnation $D(X)$ is modules over the algebra,
whereas base-changing $D\bigl((S^2)^{\mathrm{top}}_{\mathbb C_{\mathrm{gas}}}\bigr)$
to $\AnSpec(\mathbb C)$ recovers the derived category of sheaves of $\mathbb
C$-vector spaces on the underlying topological space $S^2$:
\begin{equation*}
D(X)\otimes_{D(\mathbb C_{\mathrm{gas}})}D(\mathbb C)\ \simeq\ D\bigl(\mathrm{Shv}(X,\mathbb C)\bigr).
\end{equation*}
\end{remark}

\subsection{Construction of the Tate elliptic curve}
\label{subsec:19:tate}

We return to the example that opened the discussion of analytic stacks
(\cref{lec:12,lec:14}): the Tate elliptic curve $E_q$. With the language of
analytic stacks in place, it can now be constructed properly.

\begin{remark}[The gaseous base ring]
\label{rmk:19:gaseous-base}
Recall (\cref{thm:14:gaseous-base}) the analytic ring
\begin{equation*}
A:=\mathbb Z[\widehat q^{\pm1}]_{\mathrm{gas}},
\end{equation*}
initial (as an analytic, or animated analytic, ring) among those endowed with a
topologically nilpotent \emph{unit} $q$ such that, with
$P=\mathbb Z[\mathbb N\cup\{\infty\}]/[\infty]$,
\begin{equation*}
1-q\cdot\mathrm{shift}\colon\ P\otimes_{\mathbb Z}\mathbb Z[\widehat q^{\pm1}]_{\mathrm{gas}}
\ \xrightarrow{\ \sim\ }\ P\otimes_{\mathbb Z}\mathbb Z[\widehat q^{\pm1}]_{\mathrm{gas}}
\end{equation*}
is an isomorphism. Its underlying condensed ring is the ring of Laurent series in
$q$ with integer coefficients of polynomial growth,
\begin{equation*}
\ur A(\ast)=\Bigl\{\,f=\textstyle\sum_{n} a_n q^{n}\ \Big|\ a_n\in\mathbb Z,\
a_n=0\ (n\ll0),\ \exists k:\ |a_n|=O(n^{k})\ (n\to\infty)\,\Bigr\}.
\end{equation*}
\end{remark}

\begin{remark}[Goal and outline]
\label{rmk:19:tate-goal}
The goal is to define an elliptic curve $E_q$ over the discrete ring $\ur A(\ast)$
such that, as an analytic stack over $A$, its base change is the Tate uniformisation
\begin{equation*}
(E_q)^{\mathrm{sch}}_A\ =\ \mathbb{G}_{m,A}^{\mathrm{an}}\big/\,q^{\mathbb Z}.
\end{equation*}
The outline (\cref{lec:14}) is:
\begin{itemize}
\item[\textbf{Step 1.}] Produce a ``norm'' morphism of analytic stacks
$(\mathbb P^1)^{\mathrm{sch}}_A\to[0,\infty]$ (in the precise sense of
\cref{def:19:norm} below), and define the analytic $\mathbb G_m$ as the preimage
of the open interval,
\begin{equation*}
\mathbb{G}_{m,A}^{\mathrm{an}}\ =\ |\cdot|^{-1}\bigl((0,\infty)\bigr)\ \subseteq\ (\mathbb P^1)^{\mathrm{sch}}_A .
\end{equation*}
\item[\textbf{Step 2.}] One may form the quotient
$\mathbb{G}_{m,A}^{\mathrm{an}}/q^{\mathbb Z}$ \emph{in analytic stacks}.
\item[\textbf{Step 3.}] (``Algebraisation''.) One shows the quotient is a proper
smooth curve; it therefore lies in the image of
$\operatorname{Sch}_{\ur A(\ast)}\to\AnStk_A$, i.e.\ is the analytification of an
honest elliptic curve $E_q$ over $\ur A(\ast)$.
\end{itemize}
\end{remark}

\begin{remark}[On the algebraisation step]
\label{rmk:19:algebraization}
Step~3 is an algebraisation theorem: a proper smooth curve over $A$, at least one
carrying a section, is the analytification of a scheme. Producing the ample line
bundle needed is not automatic even in genus one, and here it uses GAGA
(\cref{rmk:19:gaga}); in the present framework GAGA is far more general than
Serre's and does not by itself deliver, say, Riemann--Roch, for which one wants to
know a priori that the object is algebraic.
\end{remark}

The heart of Step~1 is the notion of a norm, which is not a norm on the
underlying set of $A$ but a much more geometric datum.

\begin{definition}[Norm on an analytic ring]
\label{def:19:norm}
Let $A\in\AnRing$. A \emph{norm} on $A$ is a morphism of analytic stacks
\begin{equation*}
|\cdot|\colon\ (\mathbb P^1)^{\mathrm{sch}}_A\ \longrightarrow\ [0,\infty],
\end{equation*}
where $[0,\infty]$ and $\mathbb P^1$ over $A$ arise from the condensed-set and
scheme functors base-changed to $A$ (one must fix a value for $|q|$; the
normalisation is $|q|=\tfrac12$), subject to:
\begin{enumerate}
\item \emph{Values at $0,1,\infty$:} the constant sections $0,1,\infty\in\mathbb
P^1$ go to the constant sections $0,1,\infty\in[0,\infty]$.
\item \emph{Topologically nilpotent $\Rightarrow|\cdot|\le1$:} the topologically
nilpotent locus of $\mathbb A^1$, i.e.\ the monomorphism
\begin{equation*}
\AnSpec\bigl(A[T^{\mathbb N\cup\{\infty\}}]/(T^{\infty}=0)\bigr)\ \longrightarrow\ \AnSpec\,A[T],
\end{equation*}
has the property that its composite to $[0,\infty]$ factors over
$[0,1]\subseteq[0,\infty]$. (This last is automatically a condition, not extra
data, since the maps involved are monomorphisms.)
\item \emph{Multiplicativity:} let $X\subseteq(\mathbb P^1)^3$ be the closure of
$\{xy=z\}\subseteq\Gm^3$, and let $\overline X\subseteq[0,\infty]^3$ be the closure
of $\{xy=z\}\subseteq(0,\infty)^3$. Then the norm-cubed carries $X$ into
$\overline X$:
\begin{equation}
\label{eq:19:multiplicativity}
\begin{tikzcd}[row sep=large, column sep=large]
X^{\mathrm{sch}}_A\arrow[r]\arrow[d, "|\cdot|^3_A"'] & \bigl[(\mathbb P^1)^3\bigr]^{\mathrm{sch}}_A\arrow[d, "|\cdot|^3"] \\
\overline X\arrow[r, hook] & [0,\infty]^3 .
\end{tikzcd}
\end{equation}
\end{enumerate}
Working with $\mathbb P^1$ rather than $\mathbb A^1$ is deliberate: even on
$\mathbb A^1$ one wants to allow functions of infinite norm, so one might as well
allow $\infty$ on the source; and taking closures in
\eqref{eq:19:multiplicativity} is what handles the degenerate products $0\cdot\infty$
cleanly. Equivalently, one may require the norm to be equivariant for inversion on
$\mathbb P^1$ and on $[0,\infty]$ and ask for honest multiplicativity on the
preimage of $[0,\infty)\subseteq\mathbb P^1$ (which lies inside $\mathbb A^1$, but
is in general a proper subset, since a point of $\mathbb A^1$ can still have
infinite norm).
\end{definition}

\begin{remark}[No triangle inequality]
\label{rmk:19:no-triangle}
There is deliberately \emph{no} additivity or triangle-inequality condition. In
particular $|2|$ need not satisfy $|2|\le2$; there is a theorem that on the
gaseous base ring $|2|$ can even be $\infty$. Concretely $|2|$ is a function on
$\AnSpec(A)$ valued in $[0,\infty]$, and it is surjective, so there is a locus
where one is forced to have $|2|=\infty$.
\end{remark}

\begin{theorem}[Existence and uniqueness of the norm]
\label{thm:19:norm-existence}
There is a unique norm on $\AnSpec\,\mathbb Z[q^{\pm1}]_{\mathrm{gas}}$ taking $q$
to $\tfrac12$. (Any value in $(0,1)$ would do equally: since no property of the
target was used beyond its order structure, one rescales $[0,\infty]$ by an
exponent.)
\end{theorem}

The construction rests on a Tannaka-duality description of maps into a locally
compact Hausdorff space, treated as an analytic stack.

\begin{lemma}[Tannaka duality for locally compact Hausdorff spaces]
\label{lem:19:tanaka}
Let $X$ be a finite-dimensional (metrizable) locally compact Hausdorff space,
regarded as an analytic stack $X^{\mathrm{cond}}$ via \cref{ex:19:condensed-anima},
and let $Y\in\AnStk$. Then
\begin{equation*}
\Homop(Y,\,X^{\mathrm{cond}})\ =\ \left\{
\begin{aligned}
&F\colon D(X,\mathbb Z)\to D(Y)\ \text{tensor functors, commuting with}\\
&\text{colimits, $D(\mathbb Z)$-linear, such that $!$-locally on $Y$,}\\
&\text{for every closed $Z\subseteq X$, } F(i_{Z\,!}\mathbb Z)\in D(Y)
\text{ is connective}
\end{aligned}
\right\}.
\end{equation*}
\end{lemma}

\begin{remark}[Reading the lemma]
\label{rmk:19:tanaka-reading}
On $X$ there are many idempotent algebras: for each closed subset $Z\subseteq X$
the constant sheaf $i_{Z\,!}\mathbb Z=\mathbb Z_Z$ (which for a closed immersion
equals the pushforward, $i_!=i_*$) is an idempotent commutative algebra in
$D(X,\mathbb Z)$, and these generate $D(X,\mathbb Z)$ under colimits. Hence a
colimit-preserving tensor functor $F$ is determined by the idempotent algebras
$F(\mathbb Z_Z)$, and the only conditions to impose are that their tensor
behaviour matches the intersection pattern $Z\cap Z'$ of closed subsets (this is
what ``tensor functor'' encodes) together with a connectivity constraint. The
connectivity of $F(\mathbb Z_Z)$ is what forces the corresponding idempotent
algebra downstairs to come from an honest ($\AnRing$-level) closed immersion; if
$X$ were profinite, closed subsets would be intersections of clopens, so the
$F(\mathbb Z_Z)$ would be retracts of the unit and $Y$ affine forces them
connective automatically, but in general the connectivity must be imposed
$!$-locally on $Y$. The mapping object is in fact a set, not merely an anima.
\end{remark}

\begin{remark}[Sketch: constructing the norm]
\label{rmk:19:norm-construction}
By \cref{lem:19:tanaka}, to give the norm
$|\cdot|\colon(\mathbb P^1)^{\mathrm{sch}}_A\to[0,\infty]$ it suffices to specify,
compatibly, the idempotent algebras corresponding to the closed subsets of
$[0,\infty]$. By the multiplicativity/inversion symmetry it is enough to give the
preimages of the intervals $[0,r]$, $r\in\mathbb R$ --- i.e.\ the loci $\{|T|\le r\}$
--- as over-convergent closed discs
\begin{equation*}
\AnSpec\,A\{T\}_r\ =\ D(0,r)^{t}\ \subseteq\ (\mathbb A^1)^{\mathrm{sch}}_A
\ \subseteq\ (\mathbb P^1)^{\mathrm{sch}}_A,
\end{equation*}
where
\begin{equation*}
A\{T\}_r\ =\ \Bigl\{\,\textstyle\sum_{n\ge0} a_n T^{n}\ \Big|\ a_n\in A,\
\exists\, r'>r:\ |a_n|\,r'^{\,n}\to0\,\Bigr\}
\ =\ \varinjlim_{t}\ A\bigl[(q^{t}T)^{\mathbb N\cup\{\infty\}}\bigr]\big/\bigl((q^{t}T)^{\infty}=0\bigr),
\end{equation*}
the colimit running over (rational) exponents $t$, with the normalisation
$|q|=\tfrac12$. The idea is to start from the free algebra on a null sequence and
multiply the basis vectors by suitable powers of $q$ to enforce the
over-convergence condition. One then checks --- and this is the one computation
genuinely using the gaseous structure --- that the tensor products of these
idempotent algebras (glued at $0$ and $\infty$) reproduce exactly the intersection
behaviour of the intervals in $[0,\infty]$, which by \cref{lem:19:tanaka} yields
the desired norm. Clausen promised a fuller account of this computation next time.
\end{remark}

\begin{remark}[Analytic spaces versus analytic stacks]
\label{rmk:19:spaces-vs-stacks}
In closing discussion, the question arose whether one can axiomatise a notion of
\emph{analytic space} inside analytic stacks --- a single umbrella definition
specialising to all the classical theories. This is exactly the intent of the
whole framework: the analytic stack \emph{is} the general notion, and one should
not be put off by the change of name from ``space'' to ``stack''. Scholze and
Clausen were never satisfied with any proposed axioms cutting out the
``non-stacky'' objects; a candidate condition floated was that the functor of
points be a colimit of affines along \emph{monomorphisms}. Merely asking that the
functor take values in sets already fails --- even affine objects $\AnSpec(A)$ are
not set-valued, since they can detect derived structure.
\end{remark}
